\begin{proposition}\hypertarget{tag-01G2}{}\label{cdp:st:D.1.4}
\StatementRecord{cdp:st:D.1.4}{proved}{01G2}{\SeriesRef{cdp:st:D.1.2}}

For any $f:Y\to C$ as in \ref{cdp:st:D.1.2} put $E=R^1f_*\OO_Y$.
The coherent Leray sequence gives
\[
0\longrightarrow H^1(C,\OO_C)\xrightarrow{f^*}H^1(Y,\OO_Y)
\xrightarrow{e}H^0(C,E)\longrightarrow0.               \tag{D2}\hypertarget{tag-019U}{}\label{cdp:eq:D2}
\]
At a smooth fibre where cohomological base change holds, the
restriction homomorphism is exactly the composite
\[
H^1(Y,\OO_Y)\xrightarrow{e}H^0(C,E)
\xrightarrow{\operatorname{ev}_c}E\otimes\C(c)
\xrightarrow{\simeq}H^1(F_c,\OO_{F_c}).               \tag{D3}\hypertarget{tag-019V}{}\label{cdp:eq:D3}
\]
By \cite[Theorem~2.2, pp.~4--5]{Viaclovsky2026Combs}, $E$ is locally free. For $C=\PP^1$, write
$E=\bigoplus_{j=1}^q\OO(a_j)$.
Then the rank of \eqref{cdp:eq:D3} at a general point is exactly
$\#\{j:a_j\geq0\}$. In particular, the nonnegative splitting
degrees claimed by CDP20 Proposition 5.5 are equivalent to its
claimed general-fibre surjectivity when $E$ has rank two and is
locally free.
\end{proposition}


\begin{proof}
\begin{enumerate}[label={$\langle1\rangle$\arabic*.},leftmargin=*]
\item 
\noindent\textbf{Step 1.} Stein factorization gives $f_*\OO_Y=\OO_C$. A curve has no coherent
cohomology in degree two, so the low-degree Leray sequence is \eqref{cdp:eq:D2}.
The edge map is induced by restriction to inverse images of open
sets in $C$. Restricting these same representatives to $F_c$ gives
\eqref{cdp:eq:D3}, which proves the asserted identification of the actual maps.
The proof of \cite[Theorem~2.2, pp.~4--5]{Viaclovsky2026Combs} establishes local freeness by resolving the fibre support,
normalizing a cyclic base change with reduced central fibre, and
taking the invariant direct image. For a line bundle
$\OO(a)$ on $\PP^1$, its evaluation has rank one at every point
for $a\geq0$, and its space of global sections is zero for $a<0$.
Applying this separately to the direct summands gives the formula.

\end{enumerate}
\end{proof}